\documentclass[10pt]{article}
\usepackage[a4paper, total={8in, 11.4in} ]{geometry}
\usepackage{graphicx}
\usepackage{amssymb}

\renewcommand{\baselinestretch}{1}

\title{}

\begin{document}
\fontsize{14}{10}\textbf{\underline{Theorems:}}
\par
\fontsize{12}{10}
\textbf{\textit{\underline{Week 1:}}}
\textbf{\underline{Theorem 1.1}} Let A and B be statements. Then A $\Rightarrow $ B is equivalent to $ \overline{\mbox{A}}$$ \vee $B. That is,
(A$\Rightarrow $B) $ \Leftarrow \Rightarrow $ ($ \overline{\mbox{A}}$ $ \vee $B).
\par
\textbf{\underline{Theorem 1.2}}
Let A and B be statements. Then A $ \Rightarrow $ B is equivalent to $ \overline{\mbox{B}}$ $ \Rightarrow $ $ \overline{\mbox{A}}$. That
 is, (A $ \Rightarrow $ B) $ \Leftarrow \Rightarrow $ ($ \overline{\mbox{A}}$$ \vee $B).
\par
\textbf{\underline{Theorem 1.3}}
Let A and B be statements. Then $\overline{\mbox{A} \vee \mbox{B}}$   is equivalent to $ \overline{\mbox{A}} \wedge \overline{\mbox{B}} $.
\par
\textbf{\underline{Theorem 2.4}}
Given any two real numbers a $<$ b, there is some r $\in $Q such that a $<$r $<$ b.
\par
\textbf{\underline{Theorem 2.5}}
The number $\sqrt{2}$ is irrational. 
\par
\textbf{\textit{\underline{Week 2:}}}
\par\textbf{\underline{Theorem 3.7 -AM-GM Inequality}}
 Let n$\in$N and x$_1$,$\dots$,x$_n$ $\geq$0. Then
($x_1 \mbox{·} x_2 \dots x_n$)$^{1/n}$ $ \leq \frac{(x_1 + \dots +x_n)}{n}$
\par\textbf{\underline{Theorem 3.8 -Cauchy-Schwarz Inequality}}
 Let n $\in$ N, x$_1$,...,x$_n$, y$_1$, ..., y$_n$ $\in$ R. Then
x$_1$y$_1$ + · · · + x$_n$y$_n$ $\leq$ \par
$\sqrt{x_1^2 + ... + x_n^2}$ 
$\sqrt{y_1^2 + ... + y_n^2}$
\par\textbf{\underline{Theorem 3.8 Triangle Inequality}}
For x,y $\in$ R,
$|\mbox{x + y}| \leq |\mbox{x}| + |y|.$
\par\textbf{\underline{Corollary 3.12 Reverse Triangle Inequality}}
$||x| - |y|| \leq |x+y|.$
in particular: $ |x| -|y| \leq |x+y|$
\par\textbf{\underline{Theorem 4.1 -Geometric Series}}
Let x $\neq$ 1 be a real number. Then for N $\in $ N,
$\sum_{n=1}^{N} x^n = \frac{x-x^{N+1}}{1-x} $ and $\sum_{n=0}^{N} x^n = \frac{1-x^{N+1}}{1-x} $
\par\textbf{\underline{Theorem 4.2}}
Every Integer n $\geq$ 2 can be written as a product of primes n = p$_1$...p$_k$
\par\textbf{\underline{\textit{Week 3:}}}
\textbf{\underline{Completeness Axiom for the real numbers:}}
Every nonempty subset of R that is bounded above has a least upper bound.
\par\textbf{\underline{Theorem 5.1}}
he Archimedean property holds. That is, for any x, y $\in$ R with x, y $>$ 0,
there is some n $\in$ N such that ny $>$ x.



\end{document}